We curate our problems from three coding competition websites: \leetcode{}, \atcoder{}, and \codeforces{}. 
These websites periodically host contests containing problems that assess the coding and problem-solving skills of participants.
The problems consist of a natural language problem statement along with example input-output examples, and the goal is to write a program that passes a set of hidden tests.
Further, thousands of participants participate, solving these problems thus ensuring that the problems are vetted for clarity and correctness.

\subsection{Data Collection} 
We have written \html{} scrapers for each of the above websites to collect problems and the corresponding metadata.
To ensure quality and consistency, we parse mathematical formulas and exclude problems with images.
We also exclude problems that are not suitable for grading by input-output examples, such as those that accept multiple correct answers or require the construction of data structures.
Besides parsing the problem descriptions, we also collect associated ground truth solutions and test cases whenever directly available.
Thus for each problem, we collect tuples of  natural language problem statement \problem{}, 
test cases \tests{}, and ground truth solution \solution{}.
Finally, we associate the contest date \contestdate{} to mark the release date of each problem and 
use the collected attributes to construct problems for our four scenarios (detailed in Section~\ref{subsec:scenario-specific} ahead). 

\vspace{1pt}
\noindent \textbf{Scrolling through time.}
As noted, we associate the contest date \contestdate{} for each problem.
The release date allows us the measure the performance of \llms{} over different time windows by filtering problems 
based on whether the problem release date falls within a time window (referred to as ``scrolling'' through time).
This is crucial for evaluating and comparing models trained at different times. 
Specifically, for a new model and the corresponding cutoff date (normalized to the release date if the 
training cutoff date is not published), 
we can measure the performance of the model on benchmark problems released after the cutoff date.
%So given a new model with cutoff data (defaulting to the release date), we can measure the performance of the model on the benchmark problems released after the cutoff date, estimating the performance of the model on (likely) uncontaminated problems\footnote{Note that these competitions are designed to be challenging and the problem setters are encouraged for creating \textit{new} problems with a mix of various concepts supporting the uncontaminated assumption.}.
We have developed a UI that allows comparing models on problems released during different time windows (shown in Figure~\ref{app:fig:scrolling}).

\vspace{1pt}
\noindent \textbf{Test collection.}
Tests are crucial for assessing the correctness of the generated outputs and are used in all four scenarios.
We collect tests available on platform websites whenever possible and use them for the benchmark.
Otherwise, following \citet{liu2023your}, we use a \llm{} (here \gptfourturbo{}) to generate tests for the problems.
A key difference between our test generation approach is that instead of generating inputs directly using the \llm{}, we construct generators that sample inputs based on the problem specifications using in context learning.
Details and examples of such input generators can be found in Section~\ref{appendix:input-generators}.
Finally, we collect a small fraction of failing tests from the platform for the more recent problems allowing more directed adversarial test collection.

\vspace{1pt}
\noindent \textbf{Problem difficulty.}
Competition programming has remained a challenge for \llms{}, with \gptfour{} achieving an average \codeforces{} rating (ELO) of 392, placing it in the bottom 5 percentile ~\citep{openai2023gpt}. 
This makes it difficult to compare \llms{}, as the variation in performance across models is low.
In \livecodebench{}, we collect problems of diverse difficulties as labeled in competition platforms, excluding problems that are rated above a certain threshold that are likely too difficult for even the best models\footnote{From our early explorations, we find {\scriptsize	\textsc{CodeForces}} problems being considerably more difficult than {\scriptsize \textsc{AtCoder}} and {\scriptsize \textsc{LeetCode}}  problems and thus focus primarily on the latter platforms.}.
% We use the problem difficulty ratings on these platforms and remove problems that are too challenging for current \llms{}.
Further, we use these ratings to classify problems as \easy{}, \medium{}, and \hard{} for more granular model comparisons.


\input{tables/benchmark-stats.tex}

\subsection{Platform Specific Curation}
\label{subsec:platform-specific}
We describe the curation process for each platform.

\vspace{2pt}
\noindent \textbf{\leetcode{}.}
We collect problems from all weekly and biweekly contests on \leetcode{} that have taken place after April'23.
For each problem, we collect the problems, public tests, and user solutions.
The platform also provides a difficulty label for each problem which we use to tag the problems as \easy{}, \medium{}, and \hard{}.
Since \leetcode{} provides a starter code for each problem, we also collect it and provide it to the \llm{} in the STDIN format.
Since the hidden tests are not directly available, we use our generator-based test input generation approach (Section~\ref{appendix:input-generators}) and also collect the auto grader failing tests for some of the recent problems.

\vspace{2pt}
\noindent \textbf{\atcoder{}.}
We collect problems from the \texttt{abc} (beginner round) contests on \atcoder{} that have taken place after April'23. 
We deliberately avoid the more challenging \texttt{arc} and \texttt{agc} contests which are designed for more advanced Olympiad participants.
The problems are assigned numeric difficulty ratings, and we exclude \texttt{abc} problems with a rating of more than $500$.
We also use these numeric ratings to tag the problems as \easy{}, \medium{}, and \hard{}.
Specifically, we use the rating brackets $[0-200)$, $[200-400)$, and $[400-500]$ to perform the classification.
\atcoder{} provides public and hidden tests for each problem which we directly use in the benchmark.

\vspace{2pt}
\noindent \textbf{\codeforces{}.}
We have collected problems from the Division 3 and Division 4 contests on \codeforces{}.
Notably, we find that even with this filter, the problems are harder than the other two platforms.
\codeforces{} also provides difficulty ratings for the problems which we use to tag the problems as \easy{}, \medium{}, and \hard{} using the rating brackets $\{800\}$, $(800-1000]$, and $(1000-1300]$ respectively.
Due to the higher difficulty, we only consider a small fraction of problems from \codeforces{} and semi-automatically construct test case generators, as they do not provide complete tests on the platform (long tests are truncated).

\noindent Table~\ref*{tab:platform_comparison} provides various statistics about the problems that we have collected for \livecodebench{}. 
%(\redo{appendix}).
%Specifically, for  \leetcode{}, we use the original program difficulty labels (\easy{}, \medium{}, and \hard{}) provided by the website which we find to be a good proxy for the difficulty of the problem for current \llms{}.
% \atcoder{} platform holds weekly contests at different difficulty levels and additionally provides a difficulty rating for the problems. 
% Here, we only consider problems from \texttt{abc} (beginner round) contests and further restrict to maximum difficulty rating of $500$.
% To tag difficulties, we use the difficulty rating brackets [0-200), [200-400), and [400-500] for classifying problems as \easy{}, \medium{}, and \hard{} respectively.
% Finally, we find problems in \codeforces{} to have a higher overall difficulty compared to the other two platforms. 
% We restrict problems in Division 3 and Division 4 contests and use corresponding (platform provided) rating to annotate difficulty. 
% Specifically, we use difficulty rating brackets [800-800], (800-1000], and (1000-1300] to classify problems as \easy{}, \medium{}, and \hard{} respectively.


%A key aspect of the benchmark is attributing a release date for each problem (dependent on the contest dates). This allows us to measure the performance of \llms{} over time and fairly compare models trained at different times by ``scrolling'' the benchmark over time. Specifically, for a new model with cutoff date (normalized to release date if explicit cutoff not available), we can measure the performance of the model on the benchmark on problems released after the cutoff date. This allows us to measure the performance of the model on uncontaminated unseen problems and fairly compare it with other models.

\subsection{Scenario-specific benchmark construction}
\label{subsec:scenario-specific}

\textbf{Code Generation and Self-Repair.}
We use the natural language problem statement as the problem statement for these scenarios. 
For \leetcode{}, as noted above, an additional starter code is provided for the functional input format.
For \atcoder{} and \codeforces{} problems, we use the standard input format (similar to ~\citet{hendrycks2021measuring}). 
The collected or generated tests are then used to evaluate the correctness of the generated programs.
Our final dataset consists of $511$ problem instances across the three platforms. 

\vspace{2pt}
\noindent \textbf{Code Execution.} 
We draw inspiration from the benchmark creation procedure used in ~\citet{gu2024cruxeval}. 
First, we collect a large pool of $\smallsim 2000$ \textit{correct, human-submitted solutions} from the \leetcode{} subset. 
%Because all programs are correct, there are no subtle bugs or tricky spots in the code. 
However, many of these programs have multiple nested loops, complex numerical computations, and a large number of execution steps. 
Therefore, we apply compile-time and run-time filters to ensure samples are reasonable, and we double-check this with a manual inspection. 
More details on the filtering criteria and statistics of the dataset can be found in Appendix~\ref{appendix:dataset-execution}. 
Our final dataset consists of $479$ samples from $85$ problems.

\vspace{2pt}
\noindent \textbf{Test Case Output Prediction.} 
We use the natural language problem statement from the \leetcode{} platform and the example test inputs to construct our test case output prediction dataset.
Since the example test inputs in the problems are reasonable test cases for humans to reason about and understand the problems, they also serve as ideal test inputs for \llms{} to process.
Our final dataset consists of $442$ problem instances from a total of $181$ \leetcode{} problems.

